\originalchapter{Quiz 2: 题目与解答}
试卷日期2012年11月15日, 五题总分100. 本章为 Alan Guth 于2012年11月21日发布的\textbf{官方解答中文整理}, 并补充原点条件、几何及微分步骤. 使用原试卷已纳入课堂勘误的版本.

\originalsection{旋转实心带电球}
\begin{exercise}
题1 (25分). 均匀实心球半径 $R$、总电荷 $Q$, 绕 $z$ 轴以角速度 $\omega$ 旋转. (a, 10分) 求磁偶极矩; (b, 5分) 求远处矢量势; (c, 10分) 求球内精确矢量势. 题面允许使用半径 $a$、面密度 $\sigma$ 的旋转球壳势: 内部 $A_\phi=\mu_0a\omega\sigma r\sin\theta/3$, 外部 $A_\phi=\mu_0a^4\omega\sigma\sin\theta/(3r^2)$.
\end{exercise}
\begin{solution}
体密度 $\rho=3Q/(4\pi R^3)$, 电流 $\jj=\rho\vct\omega\times\rr$. 磁矩积分给
$ m_z=\frac12\rho\omega\int(x^2+y^2)\dd^3x=Q\omega R^2/5$, 横向分量由对称性为零. 远势为 $\vct A=\mu_0Q\omega R^2\sin\theta\hat{\vct\phi}/(20\pi r^2)$.

球内观察点 $r<R$ 处, 把源分成半径 $r'$ 的薄壳, $\sigma=\rho\dd r'$. 对 $r'<r$ 用壳外公式, 对 $r'>r$ 用壳内公式:
\[
A_\phi=\frac{\mu_0\rho\omega\sin\theta}3
\left[\frac1{r^2}\int_0^r r'^4\dd r'+r\int_r^Rr'\dd r'\right]
=\frac{\mu_0Q\omega}{4\pi R^3}\left(\frac{rR^2}2-\frac{3r^3}{10}\right)\sin\theta.
\]
势方向为 $\hat{\vct\phi}$, 在表面与球外偶极势衔接. 分壳积分中选错内外分支会丢掉 $rR^2$ 项.
\end{solution}

\originalsection{变介电常数球}
\begin{exercise}
题2 (35分). 球内 $\varepsilon(r)=\eps(R/r)^2$, 球外为 $\eps$, 无自由电荷, 远处电场 $E_0\hat{\vct z}$. (a, 9分) 推出 $\nabla^2V+(\dd\ln\varepsilon/\dd r)\partial_rV=0$; (b, 4分) 说明可按 $V_\ell(r)P_\ell(\cos\theta)$ 展开; (c, 9分) 求内外径向方程和独立幂次解; (d, 9分) 用边界条件求势; (e, 4分) 求总诱导偶极矩.
\end{exercise}
\begin{solution}
(a) $\nabla\cdot(\varepsilon\nabla V)=0$, 展开并除以 $\varepsilon$ 即得. (b) 介质径向对称且外场轴对称, 正则解无 $\phi$ 依赖. 勒让德函数的完备性使每个球面上的势可展开, 并非先猜最后只会出现一阶项.
(c) 外部方程为 $r^2V_\ell''+2rV_\ell'-\ell(\ell+1)V_\ell=0$, 两支为 $r^\ell,r^{-\ell-1}$. 内部 $\dd\ln\varepsilon/\dd r=-2/r$, 一阶径向导数项抵消, 得 $r^2V_\ell''-\ell(\ell+1)V_\ell=0$, 两支为 $r^{\ell+1},r^{-\ell}$. 对 $\ell\ge1$ 原点有界排除负幂; 零阶两支 $1,r$ 虽都有限, 后者产生非零小球 $D$ 通量, 对应额外原点自由电荷, 也须排除.
(d) 方程及边界径向对称, 模态不混合, 远场只有 $\ell=1$. 写 $V_\mathrm{in}=Ar^2\cos\theta$, $V_\mathrm{out}=(-E_0r+B/r^2)\cos\theta$. 在 $r=R$ 介电常数相同, 连续势及 $D_n$ 给 $AR=-E_0+B/R^3$, $2AR=-E_0-2B/R^3$. 解得
\[
V_\mathrm{in}=-\frac{3E_0}{4R}r^2\cos\theta,\qquad
V_\mathrm{out}=\left(-E_0r+\frac{E_0R^3}{4r^2}\right)\cos\theta.
\]
介电常数在原点发散, 本解按题设理想模型处理; 场能在原点附近可积, 且无额外点自由电荷. (e) 外势的偶极系数给 $\pp=\pi\eps E_0R^3\hat{\vct z}$.
\end{solution}

\originalsection{一对磁偶极}
\begin{exercise}
题3 (20分). $\mm_1=m_0\hat{\vct z}$ 位于 $+(a/2)\hat{\vct z}$, $\mm_2=-m_0\hat{\vct z}$ 位于 $-(a/2)\hat{\vct z}$. (a, 10分) 求轴线上大正 $z$ 处 $\vct A$ 和 $\bb$ 的首个非零项; (b, 3分) 判断多极阶次; (c, 3分) 零尺寸极限应保持 $m_0a^n$ 中哪个 $n$; (d, 4分) 已知 $\jj_\mathrm{dip}=-\mm\times\nabla\delta^3(\rr-\rr_d)$, 求极限电流密度.
\end{exercise}
\begin{solution}
(a) 轴线上两项 $\mm\times\rr$ 均零, 所以 $\vct A=0$. 磁场不能因此判零, 因为旋度需要离轴导数. 对 $z\gg a>0$,
\[
B_z=\frac{\mu_0m_0}{2\pi}\left[(z-a/2)^{-3}-(z+a/2)^{-3}\right]
=\frac{3\mu_0m_0a}{2\pi z^4}+O(a^3/z^6).
\] 官方解答第8页 Eq.(3.6) 最后蓝框写成 $3\mu_0m_0a/(4\pi z^4)$, 与其上一行矛盾; 本稿按精确差的展开修正这个系数笔误.
(b) 总磁偶极矩为零, 首项 $r^{-4}$ 是磁四极. (c) 保持 $M_q=m_0a$ 固定, 故 $n=1$. (d) 两项电流合并为 $m_0\hat{\vct z}\times\nabla[\delta^3(\rr+a\hat{\vct z}/2)-\delta^3(\rr-a\hat{\vct z}/2)]$, 对测试函数泰勒展开得 $\jj=M_q\hat{\vct z}\times\nabla\partial_z\delta^3(\rr)$, 符号为正.
\end{solution}

\originalsection{均匀磁化圆柱与球}
\begin{exercise}
题4 (10分). 无限长圆柱半径 $R$, 轴为 $z$, 内部均匀磁化 $\vct M=M_0\hat{\vct z}$. (a, 5分) 求各处 $\vct H$; (b, 5分) 求各处 $\bb$.
\end{exercise}
\begin{solution}
侧面 $\vct K_b=M_0\hat{\vct\phi}$, 等价于无限螺线管, 因而内磁场 $\mu_0M_0\hat{\vct z}$, 外场零. 无自由电流, 且圆柱侧面 $\vct M\cdot\nn=0$, 不产生退磁场, $\vct H=0$ 在内外均成立. 不能用有限圆柱两端磁化面荷的退磁场代入无限长模型.
\end{solution}
\begin{exercise}
题5 (10分). 比较沿 $+z$ 均匀极化球与均匀磁化球的场线. 选择: (a) 内外方向皆同; (b) 内同外反; (c) 内反外同; (d) 内外皆反. 原题不要求理由.
\end{exercise}
\begin{solution}
选(c). 电球内场为 $-\vct P/(3\eps)$, 磁球内场为 $2\mu_0\vct M/3$, 方向相反; 外部均为沿 $+z$ 偶极矩的偶极场, 方向相同. 这段是中文重编的解释补充, 原题的答题要求只需选项.
\end{solution}
